%
% Annex -3 

\newpage
\section{ANNEX 3 - Reply to REA Letter} 

Reply to REA Letter Ref. Ares(2025)4081333 - 21/05/2025

\subsection*{Financial Issues}

\subsubsection*{\textcolor{blue}{GANIL}}
\begin{description}
    \item[ISSUE] TNA Costs correction/clarification:
Beneficiary 2 – GANIL: The 2008 number of TNA accesses was reported at 108 euros per unit (as per Annex 1), totaling 216,864, yet only 216,000 is indicated. Please clarify and/or make the necessary corrections.

    \item[ANSWER] The correct data is 2000. The number in the Technical Report has been updated.
\end{description}


\subsubsection*{\textcolor{blue}{JSI}}

\begin{description}
    \item[ISSUE] Missing explanation on Financial Statement adjustment:
Please indicate, explain and justify the adjustment to RP1 as you set it within the financial statement.

    \item[ANSWER] The explanation has been included in the adjustment section -> A1 employees row -> tab “use of resources” -> explanation on adjustments

For your convenience, we copy the text below:
\begin{quote}
    Work performance bonus (DU2) costs have been excluded from the cost reporting, as internal conditions of the "Jožef Stefan" Institute for their payment were not met. As a result, adjustment to RP1 was necessary.

\end{quote}
\end{description}


\subsubsection*{\textcolor{blue}{UCL}}
\begin{description}
    \item[ISSUE] Use of resources:
- Please clarify the cost of 1493.14 euros (training for Siemens PLC Programming) that is set as foreseen in Annex 1, yet not visible in the Annex 1 table.


    \item[ANSWER] “Use of resource” updated in the Financial Statement. The expense has been set as not foreseen.
\end{description}

\subsubsection*{\textcolor{blue}{CNRS}}
\begin{description}
    \item[ISSUE] TNA Costs correction/clarification required:
Beneficiary 9 – CNRS – ALTO: There were 168 accesses at a unit cost of 100, totaling 16,800; however, 48,000 is indicated. Please clarify and/or make the necessary adjustments.

    \item[ANSWER] The right number is 168. The Financial Statement has been updated.
\end{description}

\subsubsection*{\textcolor{blue}{ITAINNOVA}}
\begin{description}
    \item[ISSUE] Use of resources - Purchase costs - Equipment:
Equipment (MINISTAT 240) cost deviates from original budget set in Annex 1 (double the amount). Please explain and clarify where appropriate.

    \item[ANSWER] This justification has been included in the tab “use of resource”. 
For your convenience, we copy the text below:
\begin{quote}
    The original refrigeration unit budget proved insufficient, as the selected model couldn’t meet the required \SI{-20}{\celsius} to \SI{-30}{\celsius} range. A higher-performance unit (HUBER MINISTAT 240) was necessary to meet technical demands. Additional components—control software, refrigerant, and hoses—added further unforeseen costs. These upgrades were essential to complete the work package and ensure detector testing capability.
\end{quote}
\end{description}

\subsubsection*{\textcolor{blue}{IFIN-HH}}
\begin{description}
    \item[ISSUE] TNA costs correction/clarification required:
    
Beneficiary 14 – INFIN- HH TANDEM: With 5,450 accesses at 35 euros per unit, the correct total amount should be 190,750, whereas 62,860 is indicated. Please rectify or clarify this discrepancy.


    \item[ANSWER] Beneficiary initially chose not to declare all TA in the financial statement so as not to exceed the dedicated budget. Consistent with all other Beneficiaries, IFIN decided to declare all TA provided (5450).
    
The Financial Statement has been updated.

\end{description}

\subsubsection*{\textcolor{blue}{USE}}
\begin{description}
    \item[ISSUE] TNA costs corrections/clarification required:
    
Beneficiary 15 – USE: The use of 160 TNA UC was indicated at 35 euros per unit, as outlined in Annex 1. While the amount is correctly listed, it is categorized as VA. Please correct or provide clarification.

    \item[ANSWER] The Financial Statement has been updated moving the cost in the right row.

\end{description}

\subsubsection*{\textcolor{blue}{JYU}}
\begin{description}
    \item[ISSUE] Beneficiary 18 - Internally invoiced goods and services:
Internally invoiced goods and services set as foreseen in Annex 1, yet this does not seem to be included in Annex 1. please justify appropriately within the document (9316 euros - Travel costs)


    \item[ANSWER] They are the University of Jyväskylä's own apartments, which are rented out to TA users for accommodation. This expense can be considered in travel cost category.

\end{description}


\subsection*{Technical Issues}

\begin{description}
\item[ISSUE - Clerical Errors] On page 18 table , include the 72 PMs effort for IFJ-PAN. 
\item[\textcolor{gray}{ANSWER}] Included.

\item[ISSUE - Clerical Errors] On page 28 table , correct the 21 PMs effort for KIT.
\item[\textcolor{gray}{ANSWER}] Corrected. 

\item[ISSUE] There are mismatches on table 8 when it comes to TNAs provided in RP2, e.g for the following RIS:
ECT, GANIL-SPIRAL2, IRFU, RAPID, CLEAR, N\_TOF, TESTBEAM, SUPRATECH. The
AUs encoded in the system are 481, 2008, 48, 185, 0, 271, 2832, 146.
Please correct in the system/report or clarify accordingly.
\item[\textcolor{gray}{ANSWER}] Value aligned.

\item[ISSUE] Publication number 7 “Fission of 215Fr studied with spectroscopic methods” appears not to be an open access one. Could you please clarify and/or take corrective measures?
With regard to the field entitled “PID of deposited publication”. Please kindly take note of the statement set out in the GA “at the latest at the time of publication, a machine-readable electronic copy of the published version or the final peer-reviewed manuscript accepted for publication, is deposited in a trusted repository for scientific publications”.
I understand that this was not the case for all publications encoded in SYGMA and in particular for publication 7 (I see “False” in the system). Please kindly take corrective measures (deposit a copy of the publication in an Open Access compliant repository, negotiate with the publishing venue other means of immediate Open Access) where appropriate so as to fulfill the GA’s respective provisions and obligations. If this is an issue of recording in the system, please correct according the guidance here.

There are several publications suggested by the system (openAIRE) as linked to the project. Please import the ones relevant (if any) or disregard accordingly. You cannot leave pending publications. Please ensure that open access requirements, as set out in the GA article 17/Annex 5 are met for any publication imported. Otherwise please take corrective measures as discussed above.
\item[\textcolor{gray}{ANSWER}] The publication is
on ArXiv and details of this publication are updated directly in SyGMa, on the EC portal..

\item[ISSUE] The info for IJCLAB Supratech provided on page 33 of the Tech report (312AUs) does not match the one recorded in the system. Please clarify.
\item[\textcolor{gray}{ANSWER}] Clarification has been included at page 33 of the Technical Report: 

IJCLab/SUPRATECH The facility is actively participating to the significant European In-kind contribution to the PIP-II project of FermiLab (USA) achieving through RAs to perform the tests and full validation of SC cavities designed within the project and fabricated in Italy. The SUPRATECH maintains as well strong links with the CEA/IRFU Sygergium facilities, offering combined surface processing of accelerator components for R\&D activities. During the reference period SUPRATECH offered remote access to three projects, delivering 312 AU, corresponding to 47\% of the total AU promised in the GA. 

\item[ISSUE] The info for CEA-LFRU provided on page 33 of the Tech report (70\% AUs) does not match the one recorded in the system. Please clarify.
\item[\textcolor{gray}{ANSWER}] Clarification has been included at page 33 of the Technical Report:

CEA-IRFU/Synergium The CEA-IRFU Synergium hosts two platforms: MACHAFILM and CRYOMECH, specialized on material science at cryogenic temperatures. The Synergium cryogenic laboratories are involved in several European accelerator projects on  superconducting RF (SRF) and magnet technology with strong collaborative ties with major institutes in Europe: CERN, INFN, STFC, HZB, CNRS, and world-wide: FNAL and JLAB. Through these links MACHAFILM platform have performed well supporting thre TA projects for remote access, delivering 70\% of the AU promised in the GA. 
On the Service Improvements, in the MACHAFILM facility improvements on the magnetometry and
tunnelling spectroscopy was set up (Cernox temperature sensors), and X-ray diffraction fits analysis software was developed to extract the grain size and strain parameters of received samples. The funds were also used to purchase a dedicated software to model the behaviour of liquid Helium and mechanical properties at cryogenic temperature on the CRYOMECH Platform.

\item[ISSUE] The AUs encoded in the system for CERN CLEAR and INTC Rapid does not match the ones mentioned in the technical report on page 37. Please clarify.
\item[\textcolor{gray}{ANSWER}] the values are now aligned.

\item[ISSUE] CEA-LPA-UHI-100 and CLARA are not within expectations. Please correct the reference on page 39.
\item[\textcolor{gray}{ANSWER}] Clarification has been included at page 39 of the Technical Report: 

LPA UHI100 (CEA-LIDYL) We have plan to deliver beamtime as much as we will be able to do. We have already a proposal from an Italian Team that should be submitted. The LPA-UHI100 facility team is considering the technical feasibility of the proposal with the applicants, before submission. And we are working on a 2nd proposal which could be the following of the 1st one, depending on the results obtained. We will keep advertising about the TNA through EUROLABS, participating to international/national conferences this year 2025. We have already discussed with other communities and in particular physico –chemists and radiobiologists who want to use the extrem dose rate delivered by our laser-driven electron sources to test the response of biological material to such irradiations, in the context of FLASH radiotherapy, and to validate dosimeters developed for such high dose rate source.

1 beam time project (160 hours) for 2025, in preparation for submission next march 2025, to test a temporal diagnostic adapted to the extremely short duration of the LPA electron beams. The beamtime could be allocated at the second semester of 2025. We will try to deliver another beamtime period (160hours) in the first semester 2026, planning a maximum of 2 runs in 2026 (320hours max in 2026).

\item[ISSUE] The AUs encoded in the system for JSI i.e. 323 AUS do not match the ones mentioned in the technical report on page 51 (329). Please clarify.
\item[\textcolor{gray}{ANSWER}] Clarification has been included at page 51 of the Technical Report:

Since the start of EURO-LABS, 34 projects, totaling 366 Access Units (AU), have been approved by the User Selection Panel. Of these, 31 projects, accounting for 329 AU, were completed by the end of February 2025

\item[ISSUE] The AUs encoded in the system for IFJPAN (AIC-144) i.e. 296 AUs do not match the ones mentioned in the technical report on page 52 (332). Please clarify
\item[\textcolor{gray}{ANSWER}] The values are now aligned.  
\end{description}

\subsection*{Resources}
\begin{description}
\item[ISSUE] Please clarify why CNRS has declared 5.92 not foreseeable PMs under WP1, much more PMS than foreseen in WP2 and much less PMs than foreseen in WP3, WP5.
\item[\textcolor{gray}{ANSWER}] \leavevmode
\begin{itemize}
    \item WP1 – Additional PMs declared: The declared 5.92 PMs correspond to the time of Navin Alahari (GANIL-CNRS), who is acting as scientific co-coordinator of the project. Although this role involves contributions to WP1, his efforts were not initially foreseen or budgeted in the Grant Agreement. It was later decided to monitor and report his actual involvement in the project.
    \item WP3 and WP5 – Underconsumption of PMs: The lower effort declared is due to the involvement of a temporary staff member whose salary level has not yet been properly reflected in our system. This inconsistency is being resolved with our HR department, and the related adjustments will be made in RP3.
    \item WP2 – Overconsumption of PMs: “Regarding the overconsumption of WP2, it was preferred to declare times in order to be able to make requests for missions for staff. In addition, there is a great interest and therefore a large participation of the laboratory teams at EURO-LABS, which led them to work more than expected. Thus, The NuBall campaign has been very successful”
\end{itemize}

\item[ISSUE] Please clarify why Ben2 GaNiL has not used any person months in WP2, WP5 whatsoever. 
\item[\textcolor{gray}{ANSWER}] GANIL has no budget for personnel expenses.

There are staff effort people that declare working hours on the project on the CNRS HR platform due to the cooperation between the two beneficiaries.

It was expected from the beginning that GANIL would not employ anyone even though PM was listed in the “Staff Effort” table in the GA.  GANIL cannot employ staff. It done through CNRS.

\item[ISSUE] Please clarify why Ben8 RBI has used much more Person Months in WP4 than planned.
\item[\textcolor{gray}{ANSWER}] In the work package WP4 we are involved in task 4.2, sub-task T4.2.1, and in task T4.4. The most of our personnel cost (PM) is related to T4.4. The delivery of our milestone report MS26 was in December 2023, in the P2 period. Therefore, we had to spend significant Person Months in P2 to achieve the milestone.  In parallel, by the investment of our Institute, our Tandem Accelerator Facility has been a subject of significant upgrade, that include construction of the new accelerator building. Thansfer of the facility to the new building could take between 6 to 9 months. During that period  it should be closed for users. The deadline to finish the building is July 2025, with the facility closed for users until March 2026. This corresponds to the EURO-Labs P3 period. Therefore, in the P2 period we additionally intensified activities in T4.4, and increased our efforts (PMs), expecting that we will be forced to reduce activities (and PMs) in the P3 period.

\item[ISSUE] Please clarify why BEN 14 IFIN-HH has not used any person months in WP2 whatsoever. 
\item[\textcolor{gray}{ANSWER}] The personnel costs included in the grant agreement for IFIN-HH were intended to hire a PhD student to work on service improvements task of WP2, namely the streamlined and remote access subtask. Unfortunately, IFIN-HH could’t find a suitable candidate to work on the contribution to this subtask. Given that it is already late, IFIN-HH is considering either moving the funds to travel expenses or hiring an postdoc from within IFIN-HH to work on this for 2 months.

\item[ISSUE] Please clarify why BEN 19 UU has used much less PMs than foreseen in WP3.
\item[\textcolor{gray}{ANSWER}] UU had some procedural problems and the personnel expenses for this Reporting Period do not meet all the eligibility requirements (incurred but not declared).

\item[ISSUE] Please clarify why BEN22 UMGG has used much more Person Months in WP2 than planned.  
\item[\textcolor{gray}{ANSWER}]  The proposal describes that UMGC will not hire people equally over the full duration of the project. UMGC so far declared the actual PMs that have been used, so they have not normalised to the fraction of the project duration that has passed. UMGC will not declare more PMs than what they have in the project plan. UMGC have now internally counted 12.2 PM.
\end{description}
